\chapter{Animal Bites}

\section*{Definition}
Animal bites are wounds caused by an animal's teeth. Dogs cause many bites, while cat and human bites often cause small, deep wounds with a higher infection risk. Bites can damage skin, nerves, tendons, joints, or bone and can transmit infection, tetanus, or rabies.

\section*{When to See a Doctor}
\begin{itemize}
\item Any bite that breaks the skin, especially a puncture, crush injury, or wound of the hand, face, foot, genitals, or near a joint
\item Pain, swelling, and redness at the bite site
\item Signs of infection: purulent drainage, warmth, spreading erythema
\item Fever
\item Loss of function or sensation
\item Deep wounds involving joints, tendons, or bones
\item A bite from an unknown, stray, wild, or unavailable animal, or contact with a bat; urgent public-health advice may be needed for rabies risk
\end{itemize}

\section*{How It Develops}
Animal teeth can introduce bacteria deep into tissue. Cat bites are often narrow punctures, while dog bites may cause crushing and tearing. Common bacteria include Pasteurella, staphylococci, streptococci, and anaerobes. Rabies is a rare but fatal viral infection transmitted through saliva from an infected mammal; local public-health risk assessment determines whether post-exposure prophylaxis is needed.

\section*{Causes}
\begin{itemize}
\item Dog bites (most common, 80-90 percent of animal bites)
\item Cat bites (higher infection risk)
\item Human bites, including clenched-fist injuries (high infection risk and possible blood-borne virus exposure)
\item Rodent bites (rats, mice, hamsters)
\item Wild animal bites (raccoons, bats, foxes - rabies risk)
\end{itemize}

\section*{Related Symptoms}
\begin{itemize}
\item Cellulitis (wound infection)
\item Abscess formation
\item Tetanus
\item Rabies
\item Osteomyelitis (bone infection)
\item Septic arthritis
\end{itemize}

\section*{Medical Evaluation}
\begin{itemize}
\item Clinical assessment of wound depth, location, and contamination
\item Wound culture if infection is present
\item X-ray when a fracture, tooth fragment, or radiopaque foreign body is possible; ultrasound or other imaging may be used for deeper injury
\item Blood tests or cultures when there are systemic signs of infection, extensive injury, or immune compromise
\end{itemize}

\section*{Treatment Options}
\begin{itemize}
\item Prompt, copious irrigation, removal of foreign material, and surgical assessment or debridement when needed
\item Wound closure is individualized. Selected clean facial wounds may be closed after thorough cleaning; contaminated, puncture, hand, or delayed wounds are often left open or closed later.
\item Antibiotic prophylaxis for selected high-risk bites, such as human bites, deep or puncture wounds, hand or face wounds, wounds near joints, and bites in people who are immunocompromised or asplenic. Infected wounds require antibiotics active against aerobic and anaerobic bacteria.
\item Rabies post-exposure prophylaxis only after an urgent assessment by local public-health authorities or an appropriate clinician
\item Tetanus vaccination, with tetanus immune globulin in selected high-risk wounds when vaccination history is incomplete
\item Elevation and wound care follow-up
\end{itemize}

\section*{Self-Care and Home Management}
\begin{itemize}
\item Wash the wound thoroughly with soap and water for 15 minutes
\item Apply direct pressure if bleeding
\item Elevate the wound
\item Cover with a clean dressing after washing; do not apply caustic substances or attempt to close a deep wound at home
\item Seek medical evaluation for any bite that breaks the skin unless a clinician confirms it is trivial; return urgently for increasing redness, warmth, swelling, pus, fever, red streaks, numbness, or reduced movement
\end{itemize}

\section*{References}
\begin{thebibliography}{99}
\bibitem{animal_bites:ref1} Rhea SK, Weber DJ. ``Management of animal and human bites.'' \textit{N Engl J Med}. 2014;371(25):2418-2425.
\bibitem{animal_bites:ref2} Kannikeswaran N, Kamat D, et al. ``Mammalian bites.'' \textit{Clin Pediatr}. 2009;48(2):145-150.
\bibitem{animal_bites:ref3} Manning SE, Rupprecht CE, et al. ``Human rabies prevention --- United States, 2008.'' \textit{MMWR Recomm Rep}. 2008;57(RR-3):1-28.
\bibitem{animal_bites:ref4} Stevens DL, Bisno AL, et al. ``Practice guidelines for the diagnosis and management of skin and soft tissue infections.'' \textit{Clin Infect Dis}. 2014;59(2):e10-e52.
\bibitem{animal_bites:ref5} Oehler RL, Velez AP, et al. ``Bite-related and septic syndromes caused by cats and dogs.'' \textit{Lancet Infect Dis}. 2009;9(7):439-447.
\bibitem{animal_bites:ref6} Centers for Disease Control and Prevention. Rabies post-exposure prophylaxis guidance. Updated 2025.
\bibitem{animal_bites:ref7} National Institute for Health and Care Excellence. Human and animal bites: antimicrobial prescribing (NG184). Updated 2023.
\end{thebibliography}
